%% kapa: exam=1 round=36 year=2025 period=1 form=A subject=real_estate subject_ko=부동산학원론 seq=38 number=118
%% source: Q-Net 감정평가사 기출문제 (cst00309) artlSeq=5231525 file=raw/1008/5231525/2025년 제36회 감정평가사 1차시험 1교시 문제지 원본.pdf
%% conversion: kordoc 4.24.51 -> tools/gen1.py (auto); review=proofread-sonnet; flags=none
\begin{question}{118}
대상물건의 감정평가는 주된 방법의 감정평가방법을 적용하되, 다른 감정평가방식에 속한 하나 이상의 방법으로 산출한 시산가액과 비교해 합리성을 검토하여야 한다. 이러한 시산가액조정의 유의사항에 관한 설명으로 옳지 않은 것은?
\begin{choices}
  \item 단가와 총액의 관계는 시장증거(market evidence)를 바탕으로 분석한다.
  \item 관련 자료의 활용 적부를 검토한다.
  \item 가치형성요인의 분석에서 누락이 있었는지 검토한다.
  \item 각각의 시산가액에 가중치를 두어 조정이 가능하다.
  \item 시산가액 조정 시 공시지가기준법과 거래사례비교법은 같은 감정평가방식으로 본다.
\end{choices}
\end{question}
\answer{5}
